\chapter{Clipping evidence, tails and in-place samples}
\MARSChapterHeading{Clipping evidence and in-place samples}
\label{revision-lab:pcm}

Import the development lesson \texttt{020-pcm-mixing} and choose Continue in
Learn. Read the main chapter's sample representation, saturation rule and
clipping-mask derivation first. The complete \texttt{pcm-mix.asm} is the worked
observation source; \texttt{independent/pcm-subtract.asm} is your independent
exercise. Both are also supplied in \texttt{revision/labs/files} for use without
Learn. The arithmetic does not require an audio device.

Eight inline predictions lead through seven complete sources using ordinary
Debug and Run. Use one line breakpoint at a time and stop before the selected
instruction executes. Learn highlights the next requested source operation;
Debug's selected frame identifies the actual pause. Treat the two return values
as different quantities: \texttt{\$v0} counts processed samples and
\texttt{\$v1} counts strict clipping events. Keep disposable copies for changed
inputs and intentional faults.

\section{Predict samples and the two return values}

Begin with the unchanged sixteen-sample caller. Derive the outputs for sample
indices three through six. In particular, $30000+10000$ clips, while
$32767+0$ does not. Write each result as a signed value and a low/high byte pair.
Predict the total number of strict out-of-range sums before running.

Break at \texttt{mix\_ready}. The function has returned sixteen processed samples
in \texttt{\$v0} and six clips in \texttt{\$v1}. Inspect \texttt{pcm\_c}; the
caller has not yet copied those values to \texttt{processed} and
\texttt{clipped}, whose sentinels are still 99. Continue to
\texttt{wrap\_ready}. The first two result words now contain 16 and 6; the
wrapping comparison returns live \texttt{\$v0}=16 and \texttt{\$v1}=0, while
\texttt{wrapped\_count} is still 99. At \texttt{comparison\_ready}, all three
published words contain 16, 6 and 16.

Compare the stored saturated and wrapped results for exact sums 40000 and
$-40000$. Both versions keep two bytes per sample, but they assign different
values to an out-of-range sum. Then compare an exact endpoint. A stored endpoint
alone is insufficient evidence that clipping happened. If the output is right
but the count is wrong, keep the exact sum beside the stored sample while
investigating the counter.

The optional listening material in \texttt{packages/020-pcm-mixing/media} contains
input and saturated/wrapped output WAV files. They provide another way to compare
the arithmetic; use the numerical samples if playback is unavailable. This unit
processes arrays. Parsing and writing an audio file belongs to the later WAV
pipeline, and no real-time playback deadline is part of this function's contract.

\section{Read the mask as evidence, then change its unit}

Open \texttt{worked/mask-nine.asm}. It selects the following first nine
samples while retaining the remaining declared samples and arena sizes. The
first eight pairs are the main chapter's mask example; the ninth tests the tail
after a nonzero group count. Without Learn, prepare the same values in a
disposable copy and set \texttt{count} to nine.
\begin{verbatim}
A: 0,30000,32767,-30000,-32768,30000,-30000,32767,-20000
B: 0,2767,1,-2768,-1,10000,-10000,-32768,-20000
\end{verbatim}
Their exact sums are $0,32767,32768,-32768,-32769,40000,-40000,-1,-40000$.
Four of the first eight clip, followed by the ninth. Distinguish the two exact
endpoints from the just-outside values.

Pause at \texttt{mix\_saturated}, before \texttt{ceq.h}. Expand MSA vector
registers and compare halfword lanes in \texttt{\$w2} and \texttt{\$w3}.
These are the wrapping and saturated results. Lane two wraps to hexadecimal
\texttt{8000} but saturates to \texttt{7fff}. Lane three remains
\texttt{8000} in both. Follow the comparison and inversion to
\texttt{mix\_clipping\_mask}. In \texttt{\$w4}, only lanes 2, 4, 5 and 6
are \texttt{ffff}. The other halfwords are zero. The debugger displays raw
hexadecimal lane values; an all-one mask is evidence, not a signed audio sample
to add to the clipping count.

Step through \texttt{pcnt.d}, or continue to \texttt{mix\_population}.
Now \texttt{\$w5} has two 64-bit integer lanes: hexadecimal \texttt{10} and
\texttt{30}, decimal 16 and 48. The first doubleword contains four halfwords,
only one clipped; the second contains three clipped halfwords. Add their bit
counts and divide by sixteen. Predict \texttt{\$v1} at
\texttt{mix\_store\_block}, before the output store. It is four.

\begin{figure}[htbp]
  \centering
  \begin{tikzpicture}[alt={Two actual native debugger views for the nine-sample input. At mix_clipping_mask, register w4 has halfword masks 0000, 0000, ffff, 0000, ffff, ffff, ffff, 0000. At mix_population, register w5 has 64-bit integer lanes hexadecimal 10 and 30, decimal sixteen and forty-eight. Native boxes identify the relevant lane widths. Adding the bit counts and dividing by sixteen produces four clipped samples before the tail.}]
    \node[anchor=north west,inner sep=0,font=\small\bfseries] (atitle) at (0,0) {1. Halfword clipping mask in \texttt{\$w4}};
    \node[anchor=north west,inner sep=0] (a) at ([yshift=-2mm]atitle.south west) {\MARSRevisionCrop{revision/captures/pcm/pcm-mask.png}{0.97\linewidth}{0.23125}{0.180}{0.2625}{0.758}};
    \node[anchor=north west,inner sep=0,font=\small\bfseries] (btitle) at ([yshift=-5mm]a.south west) {2. Doubleword population counts in \texttt{\$w5}};
    \node[anchor=north west,inner sep=0] (b) at ([yshift=-2mm]btitle.south west) {\MARSRevisionCrop{revision/captures/pcm/pcm-population.png}{0.97\linewidth}{0.23125}{0.120}{0.2625}{0.778}};
    \begin{scope}[shift={(a.south west)},x={($(a.south east)-(a.south west)$)},y={($(a.north west)-(a.south west)$)}]
      \draw[black,line width=.7pt] (.038,.025) rectangle (.694,.36);
    \end{scope}
    \begin{scope}[shift={(b.south west)},x={($(b.south east)-(b.south west)$)},y={($(b.north west)-(b.south west)$)}]
      \draw[black,line width=.7pt] (.038,.013) rectangle (.596,.22);
    \end{scope}
  \end{tikzpicture}
  \caption{Actual light-theme Debug observations of the same first group. The upper box marks which halfwords clipped. The lower box marks counts of set bits in two doublewords. Hexadecimal \texttt{10} and \texttt{30} represent 16 and 48 bits; they are not counts of samples. Native LaTeX supplies the boxes and crops.}
  \label{fig:revision-pcm-native-mask}
\end{figure}


Continue through the scalar tail to \texttt{mix\_ready}. Both the ninth stored
sample and the clipping count must reflect the final excessive sum. A correct
vector mask with a final count of one would suggest that the tail discarded the
four earlier events. A count of 64 after the first group would suggest that bit
counts were mistaken for sample counts.

\section{Check an exact alias and the first unprocessed sample}

Open \texttt{worked/inplace-nine.asm}. Its first destination pointer refers to
\texttt{pcm\_b}; this is the only change from the nine-sample mask source.
Without Learn, make that pointer change in a copy. It creates an exact allowed
alias, not a shifted overlap.
Stop at \texttt{mix\_after\_block}. The first eight B samples have been replaced
by their results; both input pointers and the output pointer still identify the
current group because the pointer updates follow this label. Step through those
updates. Each pointer advances sixteen bytes, remaining count becomes one, and
\texttt{\$v1} remains four.

At \texttt{mix\_sample}, both input pointers identify their ninth original
sample. The earlier group store did not overwrite it. Load both values before
storing the saturated result. At return, the ninth B sample is $-32768$, the
following B sample remains its original value one, and source A is unchanged.
The returned clipping count is five. Inspecting the unchanged following sample
helps distinguish a correct tail from a vector overstore.

This variant continues to the ordinary wrapping call, but B now contains the
first call's saturated output. The later call therefore receives changed input
samples. Its output cannot establish a difference between saturation and
wrapping for the original A and B. Use the unchanged disjoint source for that
comparison; use this variant to investigate the exact alias in the first call.

\begin{figure}[htbp]
  \centering
  \begin{tikzpicture}[alt={An actual native Memory row at address hexadecimal 100100d0 after the nine-sample in-place call. The sixteen bytes are ff 7f 00 80 00 80 ff 7f 00 80 ff ff 00 80 01 00. Native box one marks bytes 00 80 at address 100100dc: the last selected sample has become negative 32768. Box two marks the following 01 00: the next unselected original sample remains one.}]
    \node[anchor=south west,inner sep=0] (shot) at (0,0) {\MARSRevisionCrop{revision/captures/pcm/pcm-inplace-tail.png}{0.97\linewidth}{0.0328125}{0.3366667}{0.584375}{0.6088889}};
    \begin{scope}[shift={(shot.south west)},x={($(shot.south east)-(shot.south west)$)},y={($(shot.north west)-(shot.south west)$)}]
      \draw[black,line width=.7pt] (.78,.04) rectangle (.864,.47);
      \draw[black,line width=.7pt] (.878,.04) rectangle (.96,.47);
      \node[font=\scriptsize\bfseries,anchor=south] at (.822,.52) {1};
      \node[font=\scriptsize\bfseries,anchor=south] at (.919,.52) {2};
    \end{scope}
  \end{tikzpicture}
  \caption{Actual native Memory observation at \texttt{mix\_ready}, with output exactly replacing B. Box 1 is the ninth selected sample, stored as \mbox{\texttt{00 80}} after its excessive negative sum. Box 2 is the next original sample, \mbox{\texttt{01 00}}, outside the selected output. Native annotations distinguish a changed tail from its preserved neighbor.}
  \label{fig:revision-pcm-inplace-tail}
\end{figure}


Repeat with output exactly equal to A, then with all pointers identical. For
identical inputs, each sample is doubled before saturation; it is not an automatic
zero result as it will be for the independent subtraction. Explain the access
order needed to preserve the original pair until both values have been read.
Do not extend the proof to a destination shifted into the middle of an input;
that overlap is excluded by this interface.

Use \texttt{worked/count-0.asm}, \texttt{count-7.asm} and
\texttt{count-8.asm} for the guided boundary observations. Predict complete
groups and scalar iterations before Debug. Then extend the argument to the
supplied \texttt{count-63.asm} and \texttt{count-64.asm} without a guided stop
for every iteration. Those sources change only the requested count; samples
beyond the original sixteen are zero. For every stop, identify the processed prefix, the
remaining count and the first unprocessed byte. At zero, the function returns
zero in both result registers without selected-buffer accesses. Across repeated
calls, a previous clipping count must not survive as the next initial value.

\section{Diagnose one arithmetic or boundary fault at a time}

First use \texttt{worked/signed-one.asm}. At \texttt{mix\_signed}, before
\texttt{addu}, the recovered operands are $-30000$ and $-10000$ in
\texttt{\$t0/\$t1}. The exact sum $-40000$ fits the scalar intermediate but
requires clamping for PCM16. At return the stored sample is $-32768$ and the
clip count is one.

In a copy, omit a scalar \texttt{seh}. Trace bytes \texttt{d0 8a} into the
assembled value, then the sum and comparison. Identify the first value that
stops representing $-30000$. The high byte being nonzero is not itself an error;
the missing signed interpretation is.

Restore the original and remove the mask inversion instead. Predict whether the
population count now counts equal or unequal results. Restore it, then omit the
final shift by four. Explain the unit of the returned count before comparing its
number. Restore the original again before moving the vector gate to seven
samples or reducing the output pointer's stride. In each case identify the first
access or state transition that violates the contract.

A valid alternative may load signed halfwords directly, including at odd
addresses on this Release 6 target. It may use another mask reduction or a
scalar schedule. Do not reject it merely because its temporary registers differ
from the worked function. Keep a schedule-specific Debug trace separate from a
behavioral interface check.

Run both original public suites against the complete program:
\begin{verbatim}
java -jar /path/to/cli.jar grade \
  packages/020-pcm-mixing/checks/kernel.kts \
  revision/labs/files/pcm-mix.asm
java -jar /path/to/cli.jar grade \
  packages/020-pcm-mixing/checks/program.kts \
  revision/labs/files/pcm-mix.asm
\end{verbatim}
If you measure scalar and vector work, keep the same values, count, alias and
function interval. A shorter instruction stream alone does not establish a host
elapsed-time or hardware speedup. The reduction and remainder are part of the
complete function cost.

\section{Write and test saturated subtraction independently}

Implement \texttt{subtract\_pcm16} in \texttt{pcm-subtract.asm}, keeping its
caller and data harness. Pass full A, B and output pointers in
\texttt{\$a0/\$a1/\$a2}, count zero through 64 in \texttt{\$a3}; return count
and clipped differences in \texttt{\$v0/\$v1}. Each selected slice has exactly
$2n$ bytes at arbitrary byte alignment. Input arrays are disjoint or exactly
identical; output is disjoint from both or exactly equal to either, including
all-identical pointers. Shifted partial overlaps are excluded.

Compute $A-B$ exactly in a sufficiently wide intermediate and clamp strictly
outside $[-32768,32767]$. Preserve all other caller data, saved registers, frame
pointer and stack pointer. Read only selected input/output caller bytes and write
only selected output bytes. Private constants and private stack storage are
allowed. Return through the supplied continuation, initialize the clipping count
on every call and perform no console I/O.

The supplied nine-pair example has five clipped differences. Explain why two
identical inputs yield zero output and zero clips even when their original
samples were near the endpoints. Choose a pair that distinguishes $A-B$ from
$B-A$, and an exact endpoint that distinguishes clipping from endpoint detection.
Then establish a scalar baseline before choosing a vector grouping and tail.
\texttt{subs\_s.h} and \texttt{subv.h} provide signed saturation and wrapping;
derive their comparison mask instead of reusing a counter without a proof.

\Needspace{6\baselineskip}
Run the public suite:
\begin{verbatim}
java -jar /path/to/cli.jar grade \
  revision/labs/checks/pcm-subtract.kts \
  /path/to/your/pcm-subtract.asm
\end{verbatim}
Keep the scalar result, a vector-group trace, a nine-sample exact-alias trace and
a zero-length call with your result. Explain your chosen reduction, its units and
its bounds. The checks accept equivalent implementations; your explanation and
observation establish the vector strategy separately.


Learn checks the current independent source in five groups: disjoint output,
output exactly A, output exactly B, identical inputs with disjoint output, and
all-identical pointers. The groups retain every original public case and add a
reused-call check. Its lengths are eight, one, eight and zero; returned clipping
counts must be six, zero, six and zero. A previous call's count must not survive
as the next call's initial value, even when the next call has no samples.

The disjoint group includes the six explicit arithmetic boundaries and complete
caller observations. At \texttt{subtract\_ready}, the function has returned but
the caller's result words still contain 99. At \texttt{results\_ready}, they
contain nine and five. A separate sixth assessment follows the complete caller
through silent normal exit. Passing an alias group or a return observation does
not complete those other checks.

\section{Use a failed check to choose the next observation}

Select Check this program, read its first failed value or forbidden access, and
keep a concrete explanation beside your edit. If values are right but the clip
count is wrong, compare exact differences with endpoint values and check the
reduction's units. If only an alias group fails, inspect which original operand
was overwritten before being read. These are questions to investigate, not
proof of a particular cause. Graduated hints can narrow the next observation.

All six assessments must pass for the current source, including unsaved editor
changes. Editing a worked or independent file clears its execution evidence;
Continue returns to the first activity that needs another observation. Reopening
retains source and quiz history. Keep your signed-value, alias and bounds
arguments alongside the finite test results rather than treating a saved pass
as evidence about a later edited program.
